% ACS Photonics Template
% Journal code: apchd5
% Manuscript types: Article, Letter, Perspective, Review
% Abstract: ≤250 words (Article) | opening paragraph serves as abstract (Letter)
% Figures: 3.25" (1-col) or 7" (2-col), TIFF/PDF 300 dpi min
% TOC graphic required (5.5 cm x 5 cm)

\documentclass[journal=apchd5,manuscript=article]{achemso}
% For Letter: manuscript=letter
% For Review: manuscript=review

\usepackage{amsmath,amssymb,bm}
\usepackage{graphicx}
\usepackage{xcolor}
\usepackage{hyperref}
\usepackage{siunitx}
\usepackage{booktabs}
\usepackage[version=4]{mhchem}

% Placeholder macro for missing figures
\newcommand{\placeholder}[1]{%
  \begin{center}
    \fbox{\parbox{0.9\columnwidth}{\centering\small\textit{[FIGURE: #1]}}}
  \end{center}
}

% ─── Authors ─────────────────────────────────────────────────────────────────
\author{Author One}
\affiliation{Department of Physics, University of Somewhere, City, Country}

\author{Author Two}
\affiliation{Department of Physics, University of Somewhere, City, Country}

\author{Corresponding Author}
\email{corresponding@email.com}
\affiliation{Department of Physics, University of Somewhere, City, Country}

% ─── Title ───────────────────────────────────────────────────────────────────
% No "New", "Novel", "First", "Superb", etc. in title
% No acronyms unless universally known
\title{Title of the ACS Photonics Paper}

% ─── Keywords ────────────────────────────────────────────────────────────────
% Up to 6 keywords (required for Review; optional for Article)
% \keywords{photonics, nanophotonics, plasmonics, metasurface}

\begin{document}

% ─── Abstract ────────────────────────────────────────────────────────────────
% ≤250 words; no citations, no abbreviations on first use
% Formula: Context → Problem → Approach → Key result (with numbers) → Significance
\begin{abstract}
Context sentence explaining the broader field and why it matters.
Problem sentence identifying the specific gap or challenge addressed here.
Approach sentence describing what was done and the key methodology.
Key result sentence with quantitative values (e.g., X-fold enhancement, \SI{532}{\nano\meter} excitation).
Significance sentence stating the implications and potential applications.
\end{abstract}

% ─── TOC Graphic ─────────────────────────────────────────────────────────────
\begin{tocentry}
  % Required: 5.5 cm wide × 5 cm tall, ≤ 50 words of caption
  \placeholder{TOC graphic: visual summary — schematic + key result}
\end{tocentry}

% ─── Introduction ────────────────────────────────────────────────────────────
\section{Introduction}

Broad context: why photonics/nanophotonics matters in this regime.\cite{ref1}
Specific challenge: what physical mechanism or application is poorly understood or demonstrated.\cite{ref2,ref3}
Prior art and limitations: what approaches have been tried, and where they fall short.\cite{ref4,ref5}
Our approach: ``Here we demonstrate...'' / ``In this work, we report...''\cite{ref6}
Brief outline of results: one or two sentences previewing the paper structure.

% ─── Results and Discussion ──────────────────────────────────────────────────
\section{Results and Discussion}

% ACS Photonics commonly combines Results and Discussion

\subsection{Sample Fabrication and Structural Characterization}

\begin{figure}[htbp]
  \centering
  \placeholder{Device schematic and structural characterization (SEM, AFM, XRD)}
  \caption{\textbf{Sample structure and characterization.}
           (a) Schematic illustration of the device architecture.
           (b) SEM image of the fabricated sample.
           (c) AFM topography map.
           (d) Optical photograph (scale bar: \SI{10}{\micro\meter}).}
  \label{fig:structure}
\end{figure}

\subsection{Optical Response}

\begin{figure*}[htbp]
  \centering
  \placeholder{Main optical characterization: transmission/reflection spectra, dispersion, near-field maps}
  \caption{\textbf{Optical characterization.}
           (a)--(b) Measured and simulated transmission spectra.
           (c) Photonic band structure.
           (d) Electric field distribution at resonance.}
  \label{fig:optical}
\end{figure*}

\subsection{Key Physical Mechanism or Device Performance}

\begin{figure}[htbp]
  \centering
  \placeholder{Key result figure: emission, modulation, detection, or sensing data}
  \caption{\textbf{[Key result].}
           Description of each panel (a)--(d).}
  \label{fig:main_result}
\end{figure}

% ─── Conclusion ──────────────────────────────────────────────────────────────
\section{Conclusion}

In summary, we have demonstrated...
The key finding—[X result with numbers]—establishes...
These results provide a route toward...
Future work will...

% ─── Methods / Experimental Section ─────────────────────────────────────────
\section{Experimental Section}
% ACS Photonics: "Experimental Section" placed after Conclusion

\subsection{Sample Preparation}

\subsection{Optical Measurements}

\subsection{Numerical Simulations}
% FDTD / COMSOL / DFT — describe software, mesh, boundary conditions

\subsection{Data Analysis}

% ─── Acknowledgments ─────────────────────────────────────────────────────────
\begin{acknowledgement}
The authors acknowledge...
This work was supported by [Funding Agency] under Grant No. XXX.
\end{acknowledgement}

% ─── Author Contributions ────────────────────────────────────────────────────
% Use CRediT taxonomy if required
% \section*{Author Contributions}
% A.O. and C.A. conceived the project. A.O. fabricated the samples.
% T.W. performed simulations. All authors analyzed data and wrote the manuscript.

% ─── Conflict of Interest ────────────────────────────────────────────────────
% \section*{Conflict of Interest}
% The authors declare no competing financial interest.

% ─── Supporting Information ──────────────────────────────────────────────────
\begin{suppinfo}
Figure S1: [description].
Figure S2: [description].
This material is available free of charge via the Internet at \url{http://pubs.acs.org}.
\end{suppinfo}

% ─── References ──────────────────────────────────────────────────────────────
\bibliography{refs}

\end{document}
